% Abstract. 
%
% \SVN $Revision: 840 $
% \SVN $Author: sgordon $
% \SVN $Date: 2014-05-12 13:39:14 +0700 (Mon, 12 May 2014) $
% \SVN $URL: https://sandilands.info/svn/Common/Styles/siitthesis/abstract.tex $

Assessment plays a central role in education, providing insights into 
student comprehension, learning progress, and curriculum effectiveness. 
Written exams are particularly valuable as they capture higher-order 
cognitive skills such as reasoning, problem-solving, and expressive 
ability. However, manual grading of written responses is both 
time-consuming and prone to inconsistency due to factors such as 
evaluator fatigue, subjective interpretation, and variations in 
teaching assistant training. In the Thai context, these challenges are 
amplified by the language’s complex grammar, tonal system, and 
context-dependent exceptions, which complicate evaluation and increase 
the likelihood of bias. Consequently, many educators rely heavily on 
multiple-choice examinations, which are easier to grade but 
insufficient for assessing deeper levels of understanding.

Recent advancements in Natural Language Processing (NLP), especially 
the development of Large Language Models (LLMs), provide new 
opportunities to address these challenges. LLMs offer sophisticated 
contextual analysis and semantic understanding that go beyond the 
limitations of traditional rule-based or keyword-based grading systems. 
However, their performance may still be unpredictable without 
structured guidance. This thesis proposes a hybrid AI-assisted grading 
framework that integrates the contextual reasoning capabilities of LLMs 
with the transparency and consistency of rule-based reasoning. The 
framework is designed specifically for Thai-language written exam 
questions, aiming to reduce grading workload, enhance fairness, and 
support wider adoption of open-ended assessments in Thai education.

Through experimental evaluation, the framework demonstrates its ability 
to align closely with human grading while minimizing inconsistency. 
The results highlight its potential as a practical solution for 
teachers facing heavy workloads, as well as a step toward more 
equitable and comprehensive assessment methods in Thai education. 
This research contributes both a novel methodology for exam evaluation 
and empirical insights into the broader use of AI in educational 
contexts.
\\\\
% \textbf{Keywords:} drowsiness detection, gradient vectors, displacement vectors, Infrared LEDs, face detection, headlight recognition, night-time, traffic surveillance system, 3D structure tensor.

\begin{table} [H]
 %\begin{minipage}{3cm}
 %\begin{center}
  %\begin{adjustbox}{width=1 \textwidth}
  \begin{tabular}{p{1.5cm} p{13cm}}
   %\hline\noalign{\smallskip}
  \setlength{\parindent}{-1.2ex} \textbf{Keywords}: & Exam Evaluation, Natural Language Processing (NLP), Large Language Models (LLMs), Thai Language, Rule-Based Reasoning, Educational Assessment.\\
  \end{tabular}
  %\end{adjustbox}
 %\end{center}
\end{table}
